
\textbf{Corollaire 6.14. (G. Hermann).} \emph{Pour tout @@M0@@, il existe un polyn\^ome @@M1@@ tel que pour tout corps @@M2@@ et tout id\'eal @@M3@@ de l'anneau de polyn\^omes @@M4@@, engendr\'e par des \'el\'ements de degr\'e @@M5@@, la racine @@M6@@ de @@M7@@ est engendr\'ee par des \'el\'ements de degr\'e @@M8@@.}

En effet, soit @@M9@@ la cl\^oture alg\'ebrique de @@M10@@. Introduisant une ind\'etermin\'ee auxiliaire, consid\'erons @@M11@@ (resp. @@M12@@) le sous-sch\'ema de @@M13@@ d\'efini par l'id\'eal homog\'en\'eis\'e de @@M14@@ (resp. de @@M15@@). Dans le polyn\^ome de Hilbert @@M16@@, on a @@M17@@ avec @@M18@@ d'apr\`es 6.13.

Soit @@M19@@ la d\'ecomposition suivant la dimension. Alors @@M20@@ avec @@M21@@. En effet, quitte \`a couper @@M22@@ (resp. @@M23@@) par un espace lin\'eaire ``g\'en\'eral'' de codimension @@M24@@, on peut supposer @@M25@@. Alors \'evidemment, @@M26@@ et @@M27@@, et en vertu de 6.4 (iii), @@M28@@.

Par cons\'equent, @@M29@@ est un @@M30@@-faisceau gr\^ace \`a 6.10. Il r\'esulte donc de l'injection @@M31@@ que @@M32@@ est un @@M33@@-faisceau. Par suite, dans le polyn\^ome de Hilbert @@M34@@, on a @@M35@@ avec @@M36@@ d'apr\`es 6.7. Enfin l'id\'eal @@M37@@ de @@M38@@ est @@M39@@-r\'egulier avec @@M40@@ ; d'o\`u l'assertion gr\^ace \`a 1.3.

